\documentclass[11pt,a4paper]{article}
\usepackage{../shared/luxcover}
\usepackage[utf8]{inputenc}
\usepackage[T1]{fontenc}
\usepackage{amsmath,amssymb,amsfonts,amsthm}
\usepackage{hyperref}
\usepackage{booktabs}
\usepackage{enumitem}
\usepackage{algorithm}
\usepackage{algpseudocode}
\usepackage{xcolor}
\usepackage{tikz}
\usetikzlibrary{arrows.meta,positioning,shapes.geometric,fit}

\hypersetup{colorlinks=true,linkcolor=black,citecolor=blue,urlcolor=blue}

\newtheorem{theorem}{Theorem}[section]
\newtheorem{lemma}[theorem]{Lemma}
\newtheorem{corollary}[theorem]{Corollary}
\newtheorem{proposition}[theorem]{Proposition}
\newtheorem{definition}{Definition}[section]
\newtheorem{remark}{Remark}[section]

\title{Chain-First Architecture:\\Blockchain as Source of Truth for\\Regulated Financial Infrastructure}

\author{
Lux Research\thanks{Corresponding author: research@lux.network}\\
\textit{Lux Industries, Inc.}\\
\texttt{research@lux.network}
}

\date{2024}

\begin{document}

\luxcoverpage

\begin{abstract}
We describe a ``chain-first'' architecture for regulated financial
infrastructure in which the blockchain is the authoritative source of
truth for all asset ownership and transfer records, while traditional
databases serve only as read caches that are always re-derivable from
on-chain state.  This inverts the conventional pattern in which a
centralized database is authoritative and the blockchain is an
afterthought settlement layer.  We formalize the architecture as a
state machine with two components: (1)~an append-only chain that
records every state transition via digitally signed transactions, and
(2)~a local cache (SQLite in development, PostgreSQL in production)
that materializes the current state by replaying chain events through
a deterministic indexer.  We prove three correctness properties:
\emph{cache consistency} (the cache always reflects the chain state),
\emph{crash recovery} (the cache can be fully reconstructed from the
chain), and \emph{chain independence} (the API layer continues to
serve reads from the cache when the chain is temporarily unavailable,
and queues writes for later settlement).  The architecture is deployed
on a regulated EVM L1 for an ATS handling
12,784 security tokens, with benchmarks showing 2ms read latency from
cache versus 150ms from chain RPC.
\end{abstract}

\tableofcontents
\newpage

%% =====================================================================

\section{Introduction}

\label{sec:intro}
%% =====================================================================

Financial systems have two fundamental requirements for record-keeping:
\emph{durability} (records must not be lost) and \emph{integrity}
(records must not be silently modified).  Traditional architectures
use a centralized database (PostgreSQL, Oracle, SQL Server) as the
source of truth, protected by access controls, audit logs, and backup
procedures.  The integrity guarantee is institutional: it holds because
the operator is trustworthy and the database is properly administered.

Blockchain provides a stronger guarantee.  Every state transition is
recorded as a digitally signed transaction in an append-only ledger.
The integrity guarantee is cryptographic: modifying any record requires
forging a digital signature or controlling a supermajority of
validators.  The durability guarantee is distributed: the ledger is
replicated across every full node.

The chain-first architecture makes the blockchain the single source of
truth and reduces the database to a materialized view---a read cache
that accelerates queries but contains no original data.  If the cache
is lost, it is fully reconstructable from the chain.  If the chain
and cache disagree, the chain wins.

\subsection{Why Not Database-First?}

In a database-first architecture, the blockchain is a settlement
rail: the database records a trade, and a background job eventually
posts the settlement transaction to the chain.  This creates several
problems:

\begin{enumerate}[nosep]
  \item \textbf{Split brain.}  If the database says Alice owns 100
        shares but the chain says she owns 95, which is correct?  The
        database-first architecture has no deterministic answer.
  \item \textbf{Silent corruption.}  A database administrator can
        modify records without leaving a cryptographic trace.  The
        audit log is stored in the same database it is supposed to
        protect.
  \item \textbf{Backup divergence.}  Database backups taken at
        different times may reflect inconsistent states.  Restoring
        from backup can lose or duplicate settlements.
  \item \textbf{Regulatory risk.}  SEC Rule 17a-4 requires records
        to be stored in non-rewritable, non-erasable format (WORM).
        A mutable database does not satisfy this requirement without
        additional infrastructure.
\end{enumerate}

\subsection{Chain-First: The Inversion}

In the chain-first architecture:

\begin{itemize}[nosep]
  \item The chain is WORM (Write Once, Read Many) by construction.
  \item Every state transition is a signed transaction.
  \item The database is a read cache, always derivable from chain
        events.
  \item If the cache is corrupted or lost, it is rebuilt by replaying
        the chain from genesis.
\end{itemize}

%% =====================================================================

\section{Formal Model}

\label{sec:model}
%% =====================================================================

\subsection{System State}

\begin{definition}[Chain State]
The chain state at block $h$ is a function
$\mathcal{C}_h : \text{Address} \times \text{Token} \rightarrow
\mathbb{N}$ mapping each (address, token) pair to a balance.
The genesis state $\mathcal{C}_0$ is defined by the genesis
configuration.  For $h > 0$:
\[
  \mathcal{C}_h = \mathcal{T}_h(\mathcal{C}_{h-1})
\]
where $\mathcal{T}_h$ is the deterministic state transition function
defined by the transactions in block $h$.
\end{definition}

\begin{definition}[Cache State]
The cache state $\mathcal{D}_h$ is maintained by the indexer.  For
each chain event $e$ emitted by block $h$, the indexer applies a
deterministic update $\delta(e)$ to the cache:
\[
  \mathcal{D}_h = \delta(e_1) \circ \delta(e_2) \circ \cdots \circ
  \delta(e_k)(\mathcal{D}_{h-1})
\]
where $e_1, \ldots, e_k$ are the events in block $h$.
\end{definition}

\begin{definition}[Consistency]
The system is \emph{consistent at block $h$} if and only if:
\[
  \forall \, (a, t) : \mathcal{D}_h(a, t) = \mathcal{C}_h(a, t)
\]
That is, the cache reflects the chain state for every address and
token.
\end{definition}

\subsection{Settlement Pipeline}

The pipeline from API request to on-chain settlement:

\begin{enumerate}[nosep]
  \item \textbf{API receives order.}  The order is validated
        (authentication, authorization, compliance checks) and
        recorded in the \texttt{pending\_orders} table.
  \item \textbf{Matching engine matches order.}  The matched trade
        is recorded in the \texttt{trades} table and the
        \texttt{pending\_settlements} table.
  \item \textbf{Settlement cron fires (30s interval).}  For each
        pending settlement:
        \begin{enumerate}[nosep]
          \item Construct the on-chain transaction (mint target token,
                burn source token).
          \item Initiate MPC signing (2-of-3 threshold).
          \item Broadcast the signed transaction.
        \end{enumerate}
  \item \textbf{Chain confirms transaction.}  The transaction is
        included in a block.
  \item \textbf{Indexer processes block.}  The indexer reads the
        block's events and updates the cache.  The
        \texttt{pending\_settlements} entry is marked as confirmed.
  \item \textbf{Cache serves reads.}  The API layer reads balances,
        positions, and transaction history from the cache.
\end{enumerate}

%% =====================================================================

\section{Correctness Properties}

\label{sec:correctness}
%% =====================================================================

\subsection{Cache Consistency}

\begin{theorem}[Cache Consistency]
\label{thm:consistency}
If the indexer processes every block from genesis through block $h$
without skipping or reordering, then $\mathcal{D}_h = \mathcal{C}_h$.
\end{theorem}

\begin{proof}
By induction on block height $h$.

\emph{Base case ($h = 0$).}  The cache is initialized from the genesis
configuration, which is identical to $\mathcal{C}_0$.  Therefore
$\mathcal{D}_0 = \mathcal{C}_0$.

\emph{Inductive step ($h \to h+1$).}  Assume
$\mathcal{D}_h = \mathcal{C}_h$.  Block $h+1$ produces events
$e_1, \ldots, e_k$, and the chain state transitions to
$\mathcal{C}_{h+1} = \mathcal{T}_{h+1}(\mathcal{C}_h)$.  The indexer
applies $\delta(e_1) \circ \cdots \circ \delta(e_k)$ to
$\mathcal{D}_h$.  Because the indexer's $\delta$ functions are the
deterministic projection of the chain's state transition function
$\mathcal{T}$ onto the cached dimensions (balances, positions,
transaction history), and because the events $e_i$ encode exactly the
state changes made by $\mathcal{T}_{h+1}$, we have
$\mathcal{D}_{h+1} = \mathcal{C}_{h+1}$.
\end{proof}

\subsection{Crash Recovery}

\begin{theorem}[Crash Recovery]
\label{thm:recovery}
If the cache is destroyed at any point, it can be fully reconstructed
by replaying all chain blocks from genesis.
\end{theorem}

\begin{proof}
Initialize an empty cache $\mathcal{D}'_0 = \mathcal{C}_0$ (from the
genesis configuration).  For each block $h = 1, 2, \ldots, H$ (where
$H$ is the current chain height), apply the indexer's $\delta$
functions to the events in block $h$.  By
Theorem~\ref{thm:consistency}, the reconstructed cache
$\mathcal{D}'_H = \mathcal{C}_H$.

The reconstruction is deterministic: given the same chain, the same
cache is produced regardless of when the reconstruction occurs.  The
chain is the only input; no external state is required.
\end{proof}

\subsection{Chain Independence}

\begin{theorem}[Chain Independence]
\label{thm:independence}
If the chain becomes temporarily unavailable, the API layer continues
to serve reads from the cache (reflecting the last known chain state)
and queues writes for later settlement.  No data is lost.
\end{theorem}

\begin{proof}
Reads are served from the cache $\mathcal{D}_h$ where $h$ is the last
processed block.  The cache is a local database (SQLite or PostgreSQL)
that does not depend on chain connectivity.

Writes (new orders, trades) are recorded in the
\texttt{pending\_settlements} table in the local database.  The
settlement cron retries pending settlements on a 30-second interval.
When chain connectivity is restored, the cron processes all queued
settlements.  The \texttt{pending\_settlements} table acts as a durable
write-ahead log.

No data is lost because:
\begin{enumerate}[nosep]
  \item Reads serve stale-but-consistent data (consistent as of block
        $h$).
  \item Writes are durably queued and will be settled when the chain
        returns.
  \item The chain's append-only property guarantees that blocks
        processed before the outage are still valid after the outage.
\end{enumerate}
\end{proof}

%% =====================================================================

\section{WORM Audit Trail}

\label{sec:worm}
%% =====================================================================

The blockchain is a natural WORM (Write Once, Read Many) storage
medium: once a transaction is included in a finalized block, it cannot
be modified or deleted without forging digital signatures or controlling
a supermajority of validators.

\subsection{Hash Chain Integrity}

Each block contains a hash of the previous block, forming a hash
chain from genesis to the current head:

\[
  H_h = \mathsf{Hash}(H_{h-1} \| \mathsf{TxRoot}_h \|
  \mathsf{StateRoot}_h \| \mathsf{Timestamp}_h)
\]

\begin{theorem}[Hash Chain Tamper Detection]
\label{thm:hashchain}
Any modification to any block $h' \leq h$ in the chain is detectable
by any party that knows the current block hash $H_h$.
\end{theorem}

\begin{proof}
Suppose block $h'$ is modified, producing a different hash
$H'_{h'} \neq H_{h'}$.  Block $h'+1$ includes $H_{h'}$ in its
preimage.  Modifying block $h'$ without modifying block $h'+1$ causes
the hash chain to break: $\mathsf{Hash}(H'_{h'} \| \cdots) \neq
H_{h'+1}$.  Therefore the adversary must also modify block $h'+1$,
which in turn requires modifying block $h'+2$, and so on up to block
$h$.  This cascade requires modifying every block from $h'$ to $h$,
which requires forging the digital signatures of the block producers
for each of those blocks.  Under the EUF-CMA security of the signature
scheme, this is infeasible.

Any party that stores $H_h$ (the current head hash) can verify the
integrity of any earlier block by recomputing the hash chain.  If
the recomputed $H_h$ differs from the stored $H_h$, a modification
has occurred.
\end{proof}

\subsection{Compliance with Rule 17a-4}

SEC Rule 17a-4(f) requires that records be preserved in a format that:
\begin{enumerate}[nosep]
  \item Preserves records exclusively in a non-rewritable,
        non-erasable format (WORM).
  \item Verifies automatically the quality and accuracy of the
        recording process.
  \item Serializes the original and duplicate units of storage media
        and provides a time-date for the required period of retention.
\end{enumerate}

The chain-first architecture satisfies all three requirements:
\begin{enumerate}[nosep]
  \item The blockchain is WORM by construction (hash chain + digital
        signatures).
  \item The indexer deterministically replays chain events; any
        discrepancy between chain and cache is detected automatically.
  \item Every block carries a timestamp, and the chain provides a
        total ordering of all events with cryptographic serialization
        (block height + transaction index).
\end{enumerate}

%% =====================================================================

\section{Architecture Details}

\label{sec:architecture}
%% =====================================================================

\subsection{Cache Layer: SQLite in Development, PostgreSQL in Production}

The cache is intentionally disposable.  In development, each ATS
instance uses an embedded SQLite database (via Hanzo Base) for
zero-configuration operation.  In production, multiple ATS instances
share a PostgreSQL database for concurrent access.

In both cases, the cache can be rebuilt from the chain:

\begin{verbatim}
$ atsd cache rebuild --from-genesis
  Processing block 0...
  Processing block 1...
  ...
  Processing block 847291...
  Cache rebuilt: 12,784 tokens, 847,291 blocks, 2.3M events
  Duration: 4m 12s
\end{verbatim}

\subsection{Indexer: Deterministic Event Replay}

The indexer subscribes to chain events via WebSocket and processes
them in strict block order.  For each event type, a deterministic
handler updates the cache:

\begin{table}[h]
\centering
\begin{tabular}{ll}
\toprule
\textbf{Chain Event} & \textbf{Cache Update} \\
\midrule
\texttt{Transfer(from, to, amount)} & Debit \texttt{from}, credit
  \texttt{to} \\
\texttt{Mint(to, amount)} & Credit \texttt{to} \\
\texttt{Burn(from, amount)} & Debit \texttt{from} \\
\texttt{SecurityTokenCreated(addr, symbol)} & Insert token record \\
\texttt{ComplianceCheck(addr, result)} & Update compliance status \\
\bottomrule
\end{tabular}
\caption{Event-to-cache mapping for the deterministic indexer.}
\label{tab:events}
\end{table}

The indexer maintains a \texttt{last\_processed\_block} cursor.  On
restart, it resumes from the cursor, ensuring exactly-once processing.

\subsection{Read Path (Fast)}

The API layer reads from the cache for all queries:

\begin{table}[h]
\centering
\begin{tabular}{lrr}
\toprule
\textbf{Query} & \textbf{Cache (ms)} & \textbf{Chain RPC (ms)} \\
\midrule
Get balance & 0.8 & 45 \\
List positions & 1.2 & 150 \\
Transaction history & 2.0 & 320 \\
Token metadata & 0.5 & 25 \\
\bottomrule
\end{tabular}
\caption{Read latency: cache vs.\ chain RPC.}
\label{tab:read-latency}
\end{table}

The cache provides 50--200$\times$ faster reads than chain RPC calls.
For a securities exchange serving 1,000 concurrent users, this is the
difference between a responsive UI and an unusable one.

\subsection{Write Path (Durable)}

Writes flow through the settlement pipeline:

\begin{verbatim}
API -> validation -> pending_orders (DB)
    -> matching    -> trades (DB) + pending_settlements (DB)
    -> cron (30s)  -> MPC sign -> broadcast -> chain
    -> indexer     -> cache update -> pending_settlements.confirmed
\end{verbatim}

If any step fails, the pending settlement remains in the queue and is
retried.  The settlement is considered final only when the chain
transaction is confirmed and the indexer has processed the block.

%% =====================================================================

\section{Resilience Analysis}

\label{sec:resilience}
%% =====================================================================

\begin{table}[h]
\centering
\begin{tabular}{p{3.5cm}p{4.5cm}p{4.5cm}}
\toprule
\textbf{Failure} & \textbf{Database-First} & \textbf{Chain-First} \\
\midrule
Cache/DB corruption & Data loss (DB is truth) & Rebuild from chain
  (chain is truth) \\
Chain outage & Trades continue, no settlement & Trades continue,
  settlements queued \\
Indexer crash & N/A (no indexer) & Resume from cursor, no data loss \\
Split brain & Ambiguous ownership & Chain wins (deterministic) \\
Admin tampering & Silent modification & Detected by hash chain \\
Backup divergence & Possible inconsistency & Cache is disposable;
  chain is the backup \\
\bottomrule
\end{tabular}
\caption{Failure modes: database-first vs.\ chain-first.}
\label{tab:failures}
\end{table}

The chain-first architecture has a strictly better failure model for
every scenario except chain outage, where both architectures degrade
similarly (trades continue, settlement is deferred).

%% =====================================================================

\section{Formal Verification}

\label{sec:formal}
%% =====================================================================

The hash chain integrity property is mechanized in Lean~4:

\begin{itemize}[nosep]
  \item \texttt{Crypto.HashChainIntegrity}: proof that any modification
        to any block in the chain is detectable by comparing the
        recomputed head hash against the stored head hash.
\end{itemize}

The cache consistency and crash recovery properties are modeled in
TLA+ in the \texttt{lux/proofs/tla/} directory as the
\texttt{VaultSync} specification.

%% =====================================================================

\section{Conclusion}

\label{sec:conclusion}
%% =====================================================================

The chain-first architecture aligns the technical architecture with
the regulatory requirement: the blockchain is a WORM audit trail that
satisfies Rule 17a-4, the database is a read cache that accelerates
queries, and the settlement pipeline ensures that every trade is
durably recorded on-chain.  The architecture tolerates chain outages,
cache corruption, and node failures without data loss.  The system
is deployed on the Partner EVM with 12,784 security tokens and
sub-2ms cache reads.

\begin{thebibliography}{99}
\bibitem{sec17a4}
SEC.
\emph{Rule 17a-4: Records to be Preserved by Certain Exchange Members,
Brokers and Dealers}.
17 CFR 240.17a-4.

\bibitem{finra4511}
FINRA.
\emph{Rule 4511: General Requirements (Books and Records)}.

\bibitem{nakamoto2008}
S.~Nakamoto.
\emph{Bitcoin: A Peer-to-Peer Electronic Cash System}.
2008.

\bibitem{wood2014}
G.~Wood.
\emph{Ethereum: A Secure Decentralised Generalised Transaction Ledger}.
Yellow Paper, 2014.

\bibitem{merkle1989}
R.~Merkle.
\emph{A Certified Digital Signature}.
CRYPTO 1989.

\bibitem{lamport1982}
L.~Lamport, R.~Shostak, M.~Pease.
\emph{The Byzantine Generals Problem}.
ACM TOPLAS, 4(3):382--401, 1982.
\end{thebibliography}



\end{document}
